% duct.tex -- TikZ remake of duct_schematic.pdf
% Compile with pdflatex.
%
% A rectangular duct of width w, height h and length L, drawn in oblique
% projection with the near end open.  Geometry measured from the original:
% w = 370.7, h = 169.6, L = 242.5 drawing units, receding at 16.09 degrees.
%
% PROJECTION.  y runs left along the front edge, z up, x into the picture at
% \xang above the horizontal.  x cross y = z is *not* satisfied by a naive
% reading -- see the handedness note by the axes below.
%
% SHADING.  The original gave its faces 244, 239, 236 and 233 out of 255 -- a
% 4% spread that reads as one flat tone.  Here the three visible faces take a
% proper ramp, lightest on the top and darkest on the floor, so the box reads
% as a solid.  Each tone is decoration: the 2.00pt edges identify the form.
%
% Sizing: set the class option and \ptbase to your document's point size.
\documentclass[10pt,border=2pt]{standalone}

\usepackage{tikz}
\usetikzlibrary{arrows.meta}

\newif\ifcolour
\colourtrue
\ifcolour
  \definecolor{axcol}{HTML}{B03030}      % the triad, as in coating.tex
\else
  \colorlet{axcol}{black}
\fi
% The two outer faces are translucent, so the far edges read through them and
% the back-left corner is one continuous line rather than solid-then-dashed.
% Every face is translucent and toned by its surface normal, not by whether it
% is inside or out: the two horizontal faces share one tone and the two
% vertical faces another, which is what a single overhead light would give.
% Toning by inside/out instead made the floor and the near side wall the two
% darkest areas, which weighted the lower left of the figure for no reason.
% Depth now comes from the overlaps -- two layers composite darker than one --
% and from the edges, rather than from an arbitrary interior/exterior ramp.
\newcommand{\falpha}{0.50}
\definecolor{fhoriz}{gray}{0.86}         % ceiling and floor; composite 0.93
\definecolor{fvert}{gray}{0.72}          % both side walls;   composite 0.86

% ---- scale ----------------------------------------------------------------
\newcommand{\ptbase}{10}
\pgfmathsetmacro{\gs}{\ptbase/25}
\pgfmathsetmacro{\dsz}{\ptbase/10}
\pgfmathsetmacro{\wallw} {2.00*\dsz}     % visible edges
\pgfmathsetmacro{\arrw}  {0.85*\dsz}     % every arrow
\pgfmathsetmacro{\lgap}  {3.0/\gs}       % axis labels sit 3.0pt off their tip
% Extension lines, per ISO 129-1: from the feature to just past the dimension
% line, with a visible gap at the object so they read as annotation rather
% than geometry.  Only where the dimension is offset -- w and L are; h ends on
% the floor and ceiling themselves.
\pgfmathsetmacro{\halfe} {0.5*\wallw/\gs}  % half an edge, in units
\pgfmathsetmacro{\extgap}{1.2/\gs}
\pgfmathsetmacro{\extrun}{2.0/\gs}

% ---- geometry (drawing units; 1 unit = 1pt at original size) --------------
\newcommand{\ww}{370.7}       % duct width, along y
\newcommand{\hh}{169.6}       % duct height, along z
\newcommand{\LL}{242.5}       % duct length, along x
\newcommand{\xang}{16.09}     % receding angle
\newcommand{\arm}{76.0}       % all three axis arrows

% oblique projection: P(x,y,z)
\pgfmathsetmacro{\ex}{cos(\xang)}
\pgfmathsetmacro{\ez}{sin(\xang)}
% arguments are parenthesised: a compound argument such as \dh+2*\lgap
% would otherwise bind as -(\dh)+2*\lgap and land on the wrong side
\newcommand{\pt}[3]{({(#1)*\ex-(#2)},{(#1)*\ez+(#3)})}

\begin{document}
\begin{tikzpicture}[
    x={\gs pt}, y={\gs pt},
    head/.style   ={-{Stealth[length=6.4*\dsz pt,width=4.3*\dsz pt,
                              inset=1.2*\dsz pt]}},
    dblhead/.style={{Stealth[length=6.4*\dsz pt,width=4.3*\dsz pt,
                              inset=1.2*\dsz pt]}-{Stealth[length=6.4*\dsz pt,
                              width=4.3*\dsz pt,inset=1.2*\dsz pt]}},
    % round, not miter: the shallow corners here are 16 degrees, where a
    % miter would extend 7x the half-width -- the spike from spec section 9
    edge/.style   ={black, line width=\wallw pt, line join=round},
    axarrow/.style={axcol, line width=\arrw pt, head},
    dimen/.style  ={black, line width=\arrw pt, dblhead},
    ext/.style    ={black, line width=0.60*\dsz pt},
    lbl/.style    ={inner sep=0pt, outer sep=0pt},
  ]

  % ---- the visible faces, back to front -----------------------------------
  % The duct is closed on all four sides and open only at x = 0 and x = L, so
  % through the near opening we see the inner left wall and the inner floor,
  % and only the overlap of the two openings shows the background.  The inner
  % wall's projection ends exactly where that overlap begins, so filling both
  % and drawing the outer surfaces over them needs no clipping.
  % inner left wall, at y = +w/2, seen from inside
  \fill[fvert, fill opacity=\falpha] \pt{0}{0.5*\ww}{0} -- \pt{0}{0.5*\ww}{\hh}
             -- \pt{\LL}{0.5*\ww}{\hh} -- \pt{\LL}{0.5*\ww}{0} -- cycle;
  % floor: y from -w/2 to +w/2, x from 0 to L
  \fill[fhoriz, fill opacity=\falpha] \pt{0}{-0.5*\ww}{0} -- \pt{0}{0.5*\ww}{0}
             -- \pt{\LL}{0.5*\ww}{0} -- \pt{\LL}{-0.5*\ww}{0} -- cycle;
  % the far opening's near-side edges, visible through the near opening
  % closed outlines throughout: a vertex inside one path is a miter join, but
  % where two separate paths meet, each contributes a butt cap and the pair
  % overhangs the corner.
  % drawn before the translucent outer faces, so the far-left vertical reads
  % through the ceiling as one unbroken line
  \draw[edge] \pt{0}{0.5*\ww}{0} -- \pt{\LL}{0.5*\ww}{0}
           -- \pt{\LL}{0.5*\ww}{\hh} -- \pt{0}{0.5*\ww}{\hh} -- cycle;
  \draw[edge] \pt{0}{-0.5*\ww}{0} -- \pt{\LL}{-0.5*\ww}{0}
           -- \pt{\LL}{0.5*\ww}{0} -- \pt{0}{0.5*\ww}{0} -- cycle;
  % right-hand side wall, at y = -w/2
  \fill[fvert, fill opacity=\falpha] \pt{0}{-0.5*\ww}{0} -- \pt{0}{-0.5*\ww}{\hh}
            -- \pt{\LL}{-0.5*\ww}{\hh} -- \pt{\LL}{-0.5*\ww}{0} -- cycle;
  % ceiling, seen from outside
  \fill[fhoriz, fill opacity=\falpha] \pt{0}{-0.5*\ww}{\hh} -- \pt{0}{0.5*\ww}{\hh}
           -- \pt{\LL}{0.5*\ww}{\hh} -- \pt{\LL}{-0.5*\ww}{\hh} -- cycle;


  % ---- visible edges ------------------------------------------------------
  % outer right wall
  \draw[edge] \pt{0}{-0.5*\ww}{0} -- \pt{\LL}{-0.5*\ww}{0}
           -- \pt{\LL}{-0.5*\ww}{\hh} -- \pt{0}{-0.5*\ww}{\hh} -- cycle;
  % ceiling
  \draw[edge] \pt{0}{-0.5*\ww}{\hh} -- \pt{\LL}{-0.5*\ww}{\hh}
           -- \pt{\LL}{0.5*\ww}{\hh} -- \pt{0}{0.5*\ww}{\hh} -- cycle;
  % the near opening
  \draw[edge] \pt{0}{0.5*\ww}{0} -- \pt{0}{-0.5*\ww}{0}
           -- \pt{0}{-0.5*\ww}{\hh} -- \pt{0}{0.5*\ww}{\hh} -- cycle;

  % ---- coordinate triad at the centre of the near bottom edge ------------
  % x into the picture, y to the left, z up: x cross y = z, right-handed.
  % x is the flow direction along the duct.
  % The origin sits on the bottom front edge, at its midpoint.  y therefore
  % runs along that edge -- the swallowed case from the spec -- but here the
  % triad is red against a black edge, so the two separate by hue and the
  % overlap reads as the axis lying in the plane of the wall, which it does.
  % Labels are placed in screen coordinates, each just beyond its arrowhead.
  \draw[axarrow] (0,0) -- ({\arm*\ex},{\arm*\ez});
  \draw[axarrow] (0,0) -- (-\arm,0);
  \draw[axarrow] (0,0) -- (0,\arm);
  % Spec: offset \lgap perpendicular, on the clockwise side of the direction of
  % travel.  x is the oblique case -- rotated and hanging below its own axis.
  % y points LEFT, so its clockwise side is above, and the anchor mirrors the
  % table's rightward form: south east rather than north west.  z is the
  % plain vertical case.
  \node[lbl, anchor=north west, rotate=\xang, axcol]
       at ({\arm*\ex+\lgap*cos(\xang-90)},{\arm*\ez+\lgap*sin(\xang-90)}) {$x$};
  \node[lbl, anchor=south east, axcol] at (-\arm,\lgap) {$y$};
  \node[lbl, anchor=west,       axcol] at (\lgap,\arm)  {$z$};
  % ---- dimensions ---------------------------------------------------------
  % all three labels stand \lgap = 3.0pt clear of their own dimension line,
  % measured perpendicular to it
  \pgfmathsetmacro{\dw}{22.0}
  % w drops straight down from the two front bottom corners
  \draw[ext] \pt{0}{-0.5*\ww}{-\extgap} -- \pt{0}{-0.5*\ww}{-\dw-\extrun};
  \draw[ext] \pt{0}{0.5*\ww}{-\extgap}  -- \pt{0}{0.5*\ww}{-\dw-\extrun};
  \draw[dimen] \pt{0}{-0.5*\ww}{-\dw} -- \pt{0}{0.5*\ww}{-\dw};
  \node[lbl, anchor=north] at \pt{0}{0}{-\dw-\lgap} {$w$};
  % Mirrored to the right of the triad.  On the left the bottom arrowhead's
  % tail came within 1.2pt of the floor's receding left edge; on the right
  % that edge has ended, the nearest object line is the floor's far edge
  % 19.5pt away, and most of the arrow stands on the white see-through.
  \pgfmathsetmacro{\dh}{-0.30*\ww}
  % stop at the inner faces, not the centrelines: z = 0 and z = h are the
  % centres of 2.00pt edges, so the arrowheads were burying 1.00pt in each
  \draw[dimen] \pt{0}{\dh}{\halfe} -- \pt{0}{\dh}{\hh-\halfe};
  \node[lbl, anchor=west] at \pt{0}{\dh-\lgap}{0.5*\hh} {$h$};
  % offset perpendicular to the receding edge, so it runs parallel to it
  \pgfmathsetmacro{\dl}{22.0}
  % the offset is a screen-space vector, so it is added inside \pt's own
  % coordinate expression rather than applied as a TikZ shift
  \pgfmathsetmacro{\ox}{\dl*sin(\xang)}
  \pgfmathsetmacro{\oz}{-\dl*cos(\xang)}
  % L runs perpendicular from the near and far bottom-right corners, along the
  % same screen-space normal the dimension line is offset by
  \pgfmathsetmacro{\ux}{sin(\xang)}
  \pgfmathsetmacro{\uz}{-cos(\xang)}
  \draw[ext] ({0.5*\ww+\extgap*\ux},{\extgap*\uz})
          -- ({0.5*\ww+(\dl+\extrun)*\ux},{(\dl+\extrun)*\uz});
  \draw[ext] ({\LL*\ex+0.5*\ww+\extgap*\ux},{\LL*\ez+\extgap*\uz})
          -- ({\LL*\ex+0.5*\ww+(\dl+\extrun)*\ux},{\LL*\ez+(\dl+\extrun)*\uz});
  \draw[dimen] ({-0.5*\ww*-1+\ox},{\oz})
            -- ({\LL*\ex+0.5*\ww+\ox},{\LL*\ez+\oz});
  % offset along the line's unit normal (\ez,-\ex) by \lgap, so its
  % perpendicular clearance matches w and h rather than being set by eye
  \node[lbl, anchor=north west]
       at ({0.5*\LL*\ex+0.5*\ww+\ox+\lgap*\ez},{0.5*\LL*\ez+\oz-\lgap*\ex}) {$L$};

\end{tikzpicture}
\end{document}
